% Source: https://artofproblemsolving.com/community/c3h2673667_help_on_geometry_problem
% Stable identity: aops-topic-2673667-q01
Since $AP=BP$, triangle $ABP$ is isosceles and $\angle PAB=\angle B$. Then
\[\angle PAC=90^\circ-\angle PAB=90^\circ-\angle B=\angle C,\]
so triangle $APC$ is isosceles with $PA=PC$. Hence $BP=AP=CP$. $\blacksquare$ (So $P$ is the midpoint of the hypotenuse.)

\begin{asy}
usepackage("amsmath");
size(3cm);
defaultpen(fontsize(10pt));
pair A=(0,0),B=(0,3),C=(4,0),P=(B+C)/2;
draw(A--B--C--cycle,linewidth(0.8));draw(A--P,dashed);draw((0,0.3)--(0.3,0.3)--(0.3,0));dot(P);
label("$A$",A,SW);label("$B$",B,N);label("$C$",C,SE);label("$P$",P,NE);
void tk(pair a,pair b){pair m=(a+b)/2,u=0.12*unit(b-a);draw((m-rotate(90)*u)--(m+rotate(90)*u));}
tk(A,P);tk(B,P);tk(P,C);
\end{asy}
